% ==========================================================================
% Chương 4 — Thực nghiệm và kết quả
% Chương này chứa các vị trí placeholder cho hình before/after
% của các lab thực nghiệm.
% Khi có hình export PDF từ draw.io hoặc screenshot, đặt vào figures/.
% ==========================================================================
\chapter{Thực nghiệm và kết quả}
\label{chap:experiments}

% TODO: Viết phần giới thiệu về môi trường thực nghiệm,
% cấu hình GKE Standard/Autopilot, và phạm vi kiểm tra.

\section{Kết quả đánh giá CIS Benchmark}
\label{sec:audit-results}

% TODO: Trình bày kết quả đánh giá sáu control trên GKE Standard và Autopilot.
% Dùng bảng hoặc biểu đồ tổng hợp PASS/FAIL/REVIEW.

\section{Lab ServiceAccount Hardening}
\label{sec:sa-hardening}

% TODO: Mô tả mục tiêu, các bước thực hiện và kết quả
% của lab ServiceAccount Hardening (4.1.4, 4.1.5).

\begin{figure}[htbp]
  \centering
  \placeholderfig{sa-hardening-before-after}{13cm}{7cm}{%
    So sánh cấu hình trước và sau ServiceAccount Hardening\\[4pt]
    Before: default SA, token tự động gắn cho mọi Pod\\[2pt]
    After: SA riêng, tắt auto-mount cho Pod không cần Kubernetes API}
  \caption{So sánh cấu hình trước và sau khi tăng cường bảo mật ServiceAccount}
  \label{fig:sa-hardening-before-after}
\end{figure}

\section{Lab NetworkPolicy Hardening}
\label{sec:networkpolicy-hardening}

% TODO: Mô tả mục tiêu, các bước thực hiện và kết quả
% của lab NetworkPolicy Hardening (4.3.1, 4.3.2).

\begin{figure}[htbp]
  \centering
  \placeholderfig{networkpolicy-before-after}{13cm}{7cm}{%
    So sánh lưu lượng mạng trước và sau NetworkPolicy Hardening\\[4pt]
    Before: đo và xác minh lưu lượng thực tế giữa các workload\\[2pt]
    After: default-deny + allow rules theo kết quả đo}
  \caption{So sánh lưu lượng mạng trước và sau khi áp dụng NetworkPolicy}
  \label{fig:networkpolicy-before-after}
\end{figure}
